% Options for packages loaded elsewhere
% Options for packages loaded elsewhere
\PassOptionsToPackage{unicode}{hyperref}
\PassOptionsToPackage{hyphens}{url}
\PassOptionsToPackage{dvipsnames,svgnames,x11names}{xcolor}
%
\documentclass[
  russian,
]{article}
\usepackage{xcolor}
\usepackage[top=20mm,left=20mm,right=20mm,bottom=20mm]{geometry}
\usepackage{amsmath,amssymb}
\setcounter{secnumdepth}{5}
\usepackage{iftex}
\ifPDFTeX
  \usepackage[T1]{fontenc}
  \usepackage[utf8]{inputenc}
  \usepackage{textcomp} % provide euro and other symbols
\else % if luatex or xetex
  \usepackage{unicode-math} % this also loads fontspec
  \defaultfontfeatures{Scale=MatchLowercase}
  \defaultfontfeatures[\rmfamily]{Ligatures=TeX,Scale=1}
\fi
\usepackage{lmodern}
\ifPDFTeX\else
  % xetex/luatex font selection
  \setmainfont[]{Times New Roman}
  \setsansfont[]{Arial}
  \setmonofont[]{Courier New}
\fi
% Use upquote if available, for straight quotes in verbatim environments
\IfFileExists{upquote.sty}{\usepackage{upquote}}{}
\IfFileExists{microtype.sty}{% use microtype if available
  \usepackage[]{microtype}
  \UseMicrotypeSet[protrusion]{basicmath} % disable protrusion for tt fonts
}{}
\makeatletter
\@ifundefined{KOMAClassName}{% if non-KOMA class
  \IfFileExists{parskip.sty}{%
    \usepackage{parskip}
  }{% else
    \setlength{\parindent}{0pt}
    \setlength{\parskip}{6pt plus 2pt minus 1pt}}
}{% if KOMA class
  \KOMAoptions{parskip=half}}
\makeatother
% Make \paragraph and \subparagraph free-standing
\makeatletter
\ifx\paragraph\undefined\else
  \let\oldparagraph\paragraph
  \renewcommand{\paragraph}{
    \@ifstar
      \xxxParagraphStar
      \xxxParagraphNoStar
  }
  \newcommand{\xxxParagraphStar}[1]{\oldparagraph*{#1}\mbox{}}
  \newcommand{\xxxParagraphNoStar}[1]{\oldparagraph{#1}\mbox{}}
\fi
\ifx\subparagraph\undefined\else
  \let\oldsubparagraph\subparagraph
  \renewcommand{\subparagraph}{
    \@ifstar
      \xxxSubParagraphStar
      \xxxSubParagraphNoStar
  }
  \newcommand{\xxxSubParagraphStar}[1]{\oldsubparagraph*{#1}\mbox{}}
  \newcommand{\xxxSubParagraphNoStar}[1]{\oldsubparagraph{#1}\mbox{}}
\fi
\makeatother


\usepackage{longtable,booktabs,array}
\newcounter{none} % for unnumbered tables
\usepackage{calc} % for calculating minipage widths
% Correct order of tables after \paragraph or \subparagraph
\usepackage{etoolbox}
\makeatletter
\patchcmd\longtable{\par}{\if@noskipsec\mbox{}\fi\par}{}{}
\makeatother
% Allow footnotes in longtable head/foot
\IfFileExists{footnotehyper.sty}{\usepackage{footnotehyper}}{\usepackage{footnote}}
\makesavenoteenv{longtable}
\usepackage{graphicx}
\makeatletter
\newsavebox\pandoc@box
\newcommand*\pandocbounded[1]{% scales image to fit in text height/width
  \sbox\pandoc@box{#1}%
  \Gscale@div\@tempa{\textheight}{\dimexpr\ht\pandoc@box+\dp\pandoc@box\relax}%
  \Gscale@div\@tempb{\linewidth}{\wd\pandoc@box}%
  \ifdim\@tempb\p@<\@tempa\p@\let\@tempa\@tempb\fi% select the smaller of both
  \ifdim\@tempa\p@<\p@\scalebox{\@tempa}{\usebox\pandoc@box}%
  \else\usebox{\pandoc@box}%
  \fi%
}
% Set default figure placement to htbp
\def\fps@figure{htbp}
\makeatother



\ifLuaTeX
\usepackage[bidi=basic,shorthands=off,provide=*]{babel}
\else
\usepackage[bidi=default,shorthands=off,provide=*]{babel}
\fi
\ifPDFTeX
\else
\babelfont{rm}[]{Times New Roman}
\fi
\ifLuaTeX
  \usepackage{selnolig} % disable illegal ligatures
\fi


\setlength{\emergencystretch}{3em} % prevent overfull lines

\providecommand{\tightlist}{%
  \setlength{\itemsep}{0pt}\setlength{\parskip}{0pt}}



 


\makeatletter
\@ifpackageloaded{caption}{}{\usepackage{caption}}
\AtBeginDocument{%
\ifdefined\contentsname
  \renewcommand*\contentsname{Содержание}
\else
  \newcommand\contentsname{Содержание}
\fi
\ifdefined\listfigurename
  \renewcommand*\listfigurename{Список Иллюстраций}
\else
  \newcommand\listfigurename{Список Иллюстраций}
\fi
\ifdefined\listtablename
  \renewcommand*\listtablename{Список Таблиц}
\else
  \newcommand\listtablename{Список Таблиц}
\fi
\ifdefined\figurename
  \renewcommand*\figurename{Рисунок}
\else
  \newcommand\figurename{Рисунок}
\fi
\ifdefined\tablename
  \renewcommand*\tablename{Таблица}
\else
  \newcommand\tablename{Таблица}
\fi
}
\@ifpackageloaded{float}{}{\usepackage{float}}
\floatstyle{ruled}
\@ifundefined{c@chapter}{\newfloat{codelisting}{h}{lop}}{\newfloat{codelisting}{h}{lop}[chapter]}
\floatname{codelisting}{Список}
\newcommand*\listoflistings{\listof{codelisting}{Список Каталогов}}
\makeatother
\makeatletter
\makeatother
\makeatletter
\@ifpackageloaded{caption}{}{\usepackage{caption}}
\@ifpackageloaded{subcaption}{}{\usepackage{subcaption}}
\makeatother
\usepackage{bookmark}
\IfFileExists{xurl.sty}{\usepackage{xurl}}{} % add URL line breaks if available
\urlstyle{same}
\hypersetup{
  pdftitle={Задачи по Эконометрике-2: Качество подгонки и сравнение моделей},
  pdfauthor={Н.В. Артамонов, А. Н. Лата (МГИМО МИД России)},
  pdflang={ru},
  colorlinks=true,
  linkcolor={blue},
  filecolor={Maroon},
  citecolor={Blue},
  urlcolor={Blue},
  pdfcreator={LaTeX via pandoc}}


\title{Задачи по Эконометрике-2: Качество подгонки и сравнение моделей}
\author{Н.В. Артамонов, А. Н. Лата (МГИМО МИД России)}
\date{}
\begin{document}
\maketitle

\renewcommand*\contentsname{Содержание}
{
\setcounter{tocdepth}{3}
\tableofcontents
}

\section{Качество
подгонки}\label{ux43aux430ux447ux435ux441ux442ux432ux43e-ux43fux43eux434ux433ux43eux43dux43aux438}

\subsection{labour force equation \#1
(probit)}\label{labour-force-equation-1-probit}

Для датасета \texttt{TableF5-1} рассмотрим несколько probit-регрессий:

\begin{itemize}
\tightlist
\item
  \textbf{Модель 1:}
  \(P(LFP=1) = \Phi(\beta_0 + \beta_1 WA + \beta_2 WA^2 + \beta_3 WE + \beta_4 KL6 + \beta_5 K618 + \beta_6 CIT + \beta_7 UN + \beta_8 \log(FAMINC))\)
\item
  \textbf{Модель 2:}
  \(P(LFP=1) = \Phi(\beta_0 + \beta_1 WE + \beta_2 KL6 + \beta_3 K618 + \beta_4 CIT + \beta_5 UN + \beta_6 \log(FAMINC))\)
\item
  \textbf{Модель 3:}
  \(P(LFP=1) = \Phi(\beta_0 + \beta_1 WA + \beta_2 WA^2 + \beta_3 WE + \beta_4 KL6 + \beta_5 \log(FAMINC))\)
\item
  \textbf{Модель 4:}
  \(P(LFP=1) = \Phi(\beta_0 + \beta_1 WA + \beta_2 WA^2 + \beta_3 WE + \beta_4 KL6)\)
\end{itemize}

\textbf{Результаты оценивания:}

{\def\LTcaptype{none} % do not increment counter
\begin{longtable}[]{@{}lllll@{}}
\toprule\noalign{}
& Модель 1 & Модель 2 & Модель 3 & Модель 4 \\
\midrule\noalign{}
\endhead
\bottomrule\noalign{}
\endlastfoot
WA & 0.008 & & -0.018 & -0.006 \\
& (0.070) & & (0.069) & (0.068) \\
I(WA ** 2) & -0.001 & & -0.000 & -0.000 \\
& (0.001) & & (0.001) & (0.001) \\
WE & 0.109*** & 0.124*** & 0.109*** & 0.123*** \\
& (0.024) & (0.024) & (0.024) & (0.022) \\
KL6 & -0.851*** & -0.621*** & -0.847*** & -0.855*** \\
& (0.115) & (0.098) & (0.114) & (0.115) \\
K618 & -0.063 & 0.031 & & \\
& (0.042) & (0.036) & & \\
CIT & -0.128 & -0.165 & & \\
& (0.107) & (0.106) & & \\
UN & -0.011 & -0.017 & & \\
& (0.016) & (0.016) & & \\
np.log(FAMINC) & 0.200* & 0.170* & 0.163 & \\
& (0.105) & (0.103) & (0.101) & \\
Intercept & -2.005 & -2.673*** & -1.435 & -0.281 \\
& (1.705) & (0.957) & (1.667) & (1.503) \\
\end{longtable}
}

\emph{Signif. codes:} \texttt{***} \(p<0.01\), \texttt{**} \(p<0.05\),
\texttt{*} \(p<0.1\)

\textbf{Показатели качества:}

Для каждой регрессии вычислите следующие показатели качества подгонки
модели: \(R^2_{pseudo}\), \(adj.R^2_{pseudo}\), Cox \& Snell,
Nagelkerke/Cragg \& Uhler, Efron, McKelvey \& Zavoina. \textbf{Ответ
округлите до 3-х десятичных знаков.}

Ответ:

{\def\LTcaptype{none} % do not increment counter
\begin{longtable}[]{@{}
  >{\raggedleft\arraybackslash}p{(\linewidth - 12\tabcolsep) * \real{0.0968}}
  >{\raggedleft\arraybackslash}p{(\linewidth - 12\tabcolsep) * \real{0.1398}}
  >{\raggedleft\arraybackslash}p{(\linewidth - 12\tabcolsep) * \real{0.1828}}
  >{\raggedleft\arraybackslash}p{(\linewidth - 12\tabcolsep) * \real{0.1290}}
  >{\raggedleft\arraybackslash}p{(\linewidth - 12\tabcolsep) * \real{0.1505}}
  >{\raggedleft\arraybackslash}p{(\linewidth - 12\tabcolsep) * \real{0.0968}}
  >{\raggedleft\arraybackslash}p{(\linewidth - 12\tabcolsep) * \real{0.2043}}@{}}
\toprule\noalign{}
\begin{minipage}[b]{\linewidth}\raggedleft
Model
\end{minipage} & \begin{minipage}[b]{\linewidth}\raggedleft
pseudo.R2
\end{minipage} & \begin{minipage}[b]{\linewidth}\raggedleft
adj.pseudo.R2
\end{minipage} & \begin{minipage}[b]{\linewidth}\raggedleft
CoxSnell
\end{minipage} & \begin{minipage}[b]{\linewidth}\raggedleft
Nagelkerke
\end{minipage} & \begin{minipage}[b]{\linewidth}\raggedleft
Efron
\end{minipage} & \begin{minipage}[b]{\linewidth}\raggedleft
McKelveyZavoina
\end{minipage} \\
\midrule\noalign{}
\endhead
\bottomrule\noalign{}
\endlastfoot
1 & 0.102 & 0.085 & 0.13 & 0.175 & 0.133 & 0.207 \\
2 & 0.076 & 0.062 & 0.099 & 0.133 & 0.1 & 0.158 \\
3 & 0.097 & 0.086 & 0.125 & 0.167 & 0.127 & 0.199 \\
4 & 0.095 & 0.085 & 0.122 & 0.163 & 0.123 & 0.195 \\
\end{longtable}
}

\subsection{approve equation \#1
(logit)}\label{approve-equation-1-logit}

Для датасета \texttt{loanapp} рассмотрим несколько logit-регрессий:

\begin{itemize}
\tightlist
\item
  \textbf{Модель 1:}
  \(P(approve=1) = \Lambda(\beta_0 + \beta_1 appinc + \beta_2 appinc^2 + \beta_3 mortno + \beta_4 unem + \beta_5 dep + \beta_6 male + \beta_7 married + \beta_8 yjob + \beta_9 self)\)
\item
  \textbf{Модель 2:}
  \(P(approve=1) = \Lambda(\beta_0 + \beta_1 mortno + \beta_2 unem + \beta_3 dep + \beta_4 male + \beta_5 married + \beta_6 yjob + \beta_7 self)\)
\item
  \textbf{Модель 3:}
  \(P(approve=1) = \Lambda(\beta_0 + \beta_1 appinc + \beta_2 appinc^2 + \beta_3 mortno + \beta_4 unem + \beta_5 dep + \beta_6 married)\)
\item
  \textbf{Модель 4:}
  \(P(approve=1) = \Lambda(\beta_0 + \beta_1 appinc + \beta_2 appinc^2 + \beta_3 mortno + \beta_4 dep)\)
\end{itemize}

\textbf{Результаты оценивания:}

{\def\LTcaptype{none} % do not increment counter
\begin{longtable}[]{@{}lllll@{}}
\toprule\noalign{}
& Модель 1 & Модель 2 & Модель 3 & Модель 4 \\
\midrule\noalign{}
\endhead
\bottomrule\noalign{}
\endlastfoot
appinc & 0.004* & & 0.003 & 0.004* \\
& (0.002) & & (0.002) & (0.002) \\
I(appinc ** 2) & -0.000** & & -0.000** & -0.000** \\
& (0.000) & & (0.000) & (0.000) \\
mortno & 0.707*** & 0.770*** & 0.697*** & 0.741*** \\
& (0.175) & (0.172) & (0.175) & (0.174) \\
unem & -0.050* & -0.052* & -0.061** & \\
& (0.030) & (0.029) & (0.029) & \\
dep & -0.155** & -0.168*** & -0.157** & -0.103* \\
& (0.065) & (0.064) & (0.065) & (0.061) \\
male & -0.021 & 0.022 & & \\
& (0.188) & (0.186) & & \\
married & 0.398** & 0.417** & 0.398** & \\
& (0.163) & (0.162) & (0.156) & \\
yjob & -0.010 & -0.007 & & \\
& (0.065) & (0.065) & & \\
self & -0.359* & -0.314 & & \\
& (0.200) & (0.195) & & \\
Intercept & 1.681*** & 1.860*** & 1.709*** & 1.600*** \\
& (0.228) & (0.193) & (0.207) & (0.163) \\
\end{longtable}
}

\emph{Signif. codes:} \texttt{***} \(p<0.01\), \texttt{**} \(p<0.05\),
\texttt{*} \(p<0.1\)

\textbf{Показатели качества:}

Для каждой регрессии вычислите следующие показатели качества подгонки
модели: \(R^2_{pseudo}\), \(adj.R^2_{pseudo}\), Cox \& Snell,
Nagelkerke/Cragg \& Uhler, Efron, McKelvey \& Zavoina. \textbf{Ответ
округлите до 3-х десятичных знаков.}

Ответ:

{\def\LTcaptype{none} % do not increment counter
\begin{longtable}[]{@{}
  >{\raggedleft\arraybackslash}p{(\linewidth - 12\tabcolsep) * \real{0.0968}}
  >{\raggedleft\arraybackslash}p{(\linewidth - 12\tabcolsep) * \real{0.1398}}
  >{\raggedleft\arraybackslash}p{(\linewidth - 12\tabcolsep) * \real{0.1828}}
  >{\raggedleft\arraybackslash}p{(\linewidth - 12\tabcolsep) * \real{0.1290}}
  >{\raggedleft\arraybackslash}p{(\linewidth - 12\tabcolsep) * \real{0.1505}}
  >{\raggedleft\arraybackslash}p{(\linewidth - 12\tabcolsep) * \real{0.0968}}
  >{\raggedleft\arraybackslash}p{(\linewidth - 12\tabcolsep) * \real{0.2043}}@{}}
\toprule\noalign{}
\begin{minipage}[b]{\linewidth}\raggedleft
Model
\end{minipage} & \begin{minipage}[b]{\linewidth}\raggedleft
pseudo.R2
\end{minipage} & \begin{minipage}[b]{\linewidth}\raggedleft
adj.pseudo.R2
\end{minipage} & \begin{minipage}[b]{\linewidth}\raggedleft
CoxSnell
\end{minipage} & \begin{minipage}[b]{\linewidth}\raggedleft
Nagelkerke
\end{minipage} & \begin{minipage}[b]{\linewidth}\raggedleft
Efron
\end{minipage} & \begin{minipage}[b]{\linewidth}\raggedleft
McKelveyZavoina
\end{minipage} \\
\midrule\noalign{}
\endhead
\bottomrule\noalign{}
\endlastfoot
1 & 0.033 & 0.019 & 0.024 & 0.046 & 0.027 & 0.066 \\
2 & 0.028 & 0.017 & 0.021 & 0.039 & 0.022 & 0.06 \\
3 & 0.031 & 0.021 & 0.023 & 0.043 & 0.026 & 0.062 \\
4 & 0.024 & 0.017 & 0.018 & 0.033 & 0.02 & 0.05 \\
\end{longtable}
}

\section{Сравнение
моделей}\label{ux441ux440ux430ux432ux43dux435ux43dux438ux435-ux43cux43eux434ux435ux43bux435ux439}

\subsection{labour force equation \#1
(probit)}\label{labour-force-equation-1-probit-1}

Для датасета \texttt{TableF5-1} рассмотрим те же спецификации
probit-регрессий:

\begin{itemize}
\tightlist
\item
  \textbf{Модель 1:}
  \(P(LFP=1) = \Phi(\beta_0 + \beta_1 WA + \beta_2 WA^2 + \beta_3 WE + \beta_4 KL6 + \beta_5 K618 + \beta_6 CIT + \beta_7 UN + \beta_8 \log(FAMINC))\)
\item
  \textbf{Модель 2:}
  \(P(LFP=1) = \Phi(\beta_0 + \beta_1 WE + \beta_2 KL6 + \beta_3 K618 + \beta_4 CIT + \beta_5 UN + \beta_6 \log(FAMINC))\)
\item
  \textbf{Модель 3:}
  \(P(LFP=1) = \Phi(\beta_0 + \beta_1 WA + \beta_2 WA^2 + \beta_3 WE + \beta_4 KL6 + \beta_5 \log(FAMINC))\)
\item
  \textbf{Модель 4:}
  \(P(LFP=1) = \Phi(\beta_0 + \beta_1 WA + \beta_2 WA^2 + \beta_3 WE + \beta_4 KL6)\)
\end{itemize}

Для каждой модели вычислите показатели информационных критериев AIC \&
BIC и \(adj.R^2_{pseudo}\). \textbf{Ответ округлите до 3-х десятичных
знаков.}

Ответ:

{\def\LTcaptype{none} % do not increment counter
\begin{longtable}[]{@{}lllll@{}}
\toprule\noalign{}
& Модель 1 & Модель 2 & Модель 3 & Модель 4 \\
\midrule\noalign{}
\endhead
\bottomrule\noalign{}
\endlastfoot
WA & 0.008 & & -0.018 & -0.006 \\
& (0.070) & & (0.069) & (0.068) \\
I(WA ** 2) & -0.001 & & -0.000 & -0.000 \\
& (0.001) & & (0.001) & (0.001) \\
WE & 0.109*** & 0.124*** & 0.109*** & 0.123*** \\
& (0.024) & (0.024) & (0.024) & (0.022) \\
KL6 & -0.851*** & -0.621*** & -0.847*** & -0.855*** \\
& (0.115) & (0.098) & (0.114) & (0.115) \\
K618 & -0.063 & 0.031 & & \\
& (0.042) & (0.036) & & \\
CIT & -0.128 & -0.165 & & \\
& (0.107) & (0.106) & & \\
UN & -0.011 & -0.017 & & \\
& (0.016) & (0.016) & & \\
np.log(FAMINC) & 0.200* & 0.170* & 0.163 & \\
& (0.105) & (0.103) & (0.101) & \\
Intercept & -2.005 & -2.673*** & -1.435 & -0.281 \\
& (1.705) & (0.957) & (1.667) & (1.503) \\
AIC & 942.680 & 965.398 & 941.392 & 941.981 \\
BIC & 984.297 & 997.767 & 969.137 & 965.102 \\
Nobs & 753 & 753 & 753 & 753 \\
R2\_adj & 0.085 & 0.062 & 0.086 & 0.085 \\
\end{longtable}
}

\emph{Signif. codes:} \texttt{***} \(p<0.01\), \texttt{**} \(p<0.05\),
\texttt{*} \(p<0.1\)

Какая модель предпочтительней по информационным критериям и
\(adj.R^2_{pseudo}\)?

Ответ:

{\def\LTcaptype{none} % do not increment counter
\begin{longtable}[]{@{}lr@{}}
\toprule\noalign{}
Критерий & Регрессия \\
\midrule\noalign{}
\endhead
\bottomrule\noalign{}
\endlastfoot
AIC & 3 \\
BIC & 4 \\
adj.pseudo.R2 & 3 \\
\end{longtable}
}

\subsection{approve equation \#1
(logit)}\label{approve-equation-1-logit-1}

Для датасета \texttt{loanapp} рассмотрим несколько logit-регрессий:

\begin{itemize}
\tightlist
\item
  \textbf{Модель 1:}
  \(P(approve=1) = \Lambda(\beta_0 + \beta_1 appinc + \beta_2 appinc^2 + \beta_3 mortno + \beta_4 unem + \beta_5 dep + \beta_6 male + \beta_7 married + \beta_8 yjob + \beta_9 self)\)
\item
  \textbf{Модель 2:}
  \(P(approve=1) = \Lambda(\beta_0 + \beta_1 mortno + \beta_2 unem + \beta_3 dep + \beta_4 male + \beta_5 married + \beta_6 yjob + \beta_7 self)\)
\item
  \textbf{Модель 3:}
  \(P(approve=1) = \Lambda(\beta_0 + \beta_1 appinc + \beta_2 appinc^2 + \beta_3 mortno + \beta_4 unem + \beta_5 dep + \beta_6 married)\)
\item
  \textbf{Модель 4:}
  \(P(approve=1) = \Lambda(\beta_0 + \beta_1 appinc + \beta_2 appinc^2 + \beta_3 mortno + \beta_4 dep)\)
\end{itemize}

Для каждой модели вычислите показатели информационных критериев AIC \&
BIC и \(adj.R^2_{pseudo}\). \textbf{Ответ округлите до 3-х десятичных
знаков.}

Ответ:

{\def\LTcaptype{none} % do not increment counter
\begin{longtable}[]{@{}lllll@{}}
\toprule\noalign{}
& Модель 1 & Модель 2 & Модель 3 & Модель 4 \\
\midrule\noalign{}
\endhead
\bottomrule\noalign{}
\endlastfoot
appinc & 0.004* & & 0.003 & 0.004* \\
& (0.002) & & (0.002) & (0.002) \\
I(appinc ** 2) & -0.000** & & -0.000** & -0.000** \\
& (0.000) & & (0.000) & (0.000) \\
mortno & 0.707*** & 0.770*** & 0.697*** & 0.741*** \\
& (0.175) & (0.172) & (0.175) & (0.174) \\
unem & -0.050* & -0.052* & -0.061** & \\
& (0.030) & (0.029) & (0.029) & \\
dep & -0.155** & -0.168*** & -0.157** & -0.103* \\
& (0.065) & (0.064) & (0.065) & (0.061) \\
male & -0.021 & 0.022 & & \\
& (0.188) & (0.186) & & \\
married & 0.398** & 0.417** & 0.398** & \\
& (0.163) & (0.162) & (0.156) & \\
yjob & -0.010 & -0.007 & & \\
& (0.065) & (0.065) & & \\
self & -0.359* & -0.314 & & \\
& (0.200) & (0.195) & & \\
Intercept & 1.681*** & 1.860*** & 1.709*** & 1.600*** \\
& (0.228) & (0.193) & (0.207) & (0.163) \\
AIC & 1447.463 & 1450.862 & 1444.586 & 1451.140 \\
BIC & 1503.326 & 1495.552 & 1483.690 & 1479.072 \\
Nobs & 1971 & 1971 & 1971 & 1971 \\
R2\_adj & 0.019 & 0.017 & 0.021 & 0.017 \\
\end{longtable}
}

\emph{Signif. codes:} \texttt{***} \(p<0.01\), \texttt{**} \(p<0.05\),
\texttt{*} \(p<0.1\)

Какая модель предпочтительней по информационным критериям и
\(adj.R^2_{pseudo}\)?

Ответ:

{\def\LTcaptype{none} % do not increment counter
\begin{longtable}[]{@{}lr@{}}
\toprule\noalign{}
Критерий & Регрессия \\
\midrule\noalign{}
\endhead
\bottomrule\noalign{}
\endlastfoot
AIC & 3 \\
BIC & 4 \\
adj.pseudo.R2 & 3 \\
\end{longtable}
}




\end{document}
